% !TeX root = ../se200_identity_regimes.tex

% =============================================================================
\section{Limits and Boundary Conditions}
\label{sec:limits}
% =============================================================================

The result of this paper is deliberately narrow.
It applies only to neutral accountability substrates satisfying the
assumptions of Section~\ref{sec:formal_setting}.
It establishes a necessary lower bound on identity regimes:
a floor, not a ceiling.
It does not certify that a deployed system has chosen the correct regime for every real-world referent.

% -----------------------------------------------------------------------------
\subsection{Conditional Scope}
\label{subsec:limits_conditional_scope}
% -----------------------------------------------------------------------------

The lower bound depends on the neutral-substrate constraint developed in
\emph{Neutral Substrates}~\citep{case2026neutral} and on the formal setting fixed in this paper.
In particular, it assumes persistent disagreement, strong stable reference, minimality,
and exclusion of hidden regimes.
If those assumptions are rejected, the bound need not apply.

The result is structural rather than domain-specific.
It concerns identity behavior under transformation.
It does not depend on any particular legal theory,
policy framework, institutional vocabulary, or causal model.
It says that, under the stated assumptions, fewer than nine core regimes
cannot preserve stable reference for the carrier-basis pairs analyzed here
without either changing identity behavior or hiding a regime distinction elsewhere.

% -----------------------------------------------------------------------------
\subsection{No Upper Bound or Completeness Claim}
\label{subsec:limits_no_completeness}
% -----------------------------------------------------------------------------

The nine-regime lower-bound core is a lower bound only.
It leaves three questions open:
whether these nine regimes exhaust the coherent identity regimes;
whether the transformation basis $\Fset$ is complete for all possible record systems;
and whether some domains require additional reference kinds.

The monotonicity result of Corollary~\ref{se200.cor.Monotonicity} is one-directional.
Adding transformations cannot collapse distinctions already witnessed by $\Fset$.
A richer basis may leave the bound unchanged or force additional splits.
It cannot make the nine-regime lower bound disappear.

The inventory of Section~\ref{sec:inventory} is likewise scoped.
It identifies reference kinds required for the accountability settings considered here.
Additional domains may require additional carriers.
Additional carriers may require additional regimes.
They do not remove the regimes already forced for the carrier-basis pairs analyzed in this paper.

The paper does not claim that obligation-bearers, occurrences, or records
admit no further admissible identity bases.
For example, a record may be tracked by descriptive content or by a
particular carrier, file, signature, or chain-of-custody token.
An occurrence may be tracked as a fine-grained realization or as a coarser event complex.
An obligation-bearer may be tracked as a legal person, office-holder,
delegated role, or responsible organizational unit.
If such bases are required by a deployment, they may force additional regimes.
That possibility does not weaken the lower bound.
It confirms that the present result is a lower-bound core, not a completeness theorem.

% -----------------------------------------------------------------------------
\subsection{Excluded Systems}
\label{subsec:limits_excluded_systems}
% -----------------------------------------------------------------------------

Several classes of systems fall outside the scope of the result.

\paragraph{Interpretively dependent identity.}
A system is outside the scope if the identity condition of a referent depends
on causal, normative, evaluative, or methodological commitments supplied by an extension.
Such a system may be useful within a coordinated framework.
It is not a neutral substrate in the sense analyzed here.

\paragraph{Relation-determined identity.}
A system is outside the scope if relations introduce, modify, or
determine the identity conditions of their relata.
Relations may connect already-individuated referents.
They may not supply the identity basis that makes those referents the same or different.

\paragraph{Extension-variant identity.}
A system is outside the scope if identity conditions vary when admissible extensions are added.
Such systems violate the stable-reference requirement.
They may support local or negotiated identity.
They do not support the strong form of stable reference studied here.

\paragraph{Embedded interpretive commitments.}
A system is outside the scope if it makes substrate-level causal or normative commitments
whose validity depends on a particular interpretive framework.
The substrate may record attributed causal or normative claims.
It may not assert them as object-level substrate facts.

\paragraph{Undeclared hidden regimes.}
A system is outside the scope if it collapses regimes
and then recovers the lost distinction through roles, flags,
contextual predicates, workflow states, schema conventions, or application-level discriminators.
Such a system has not reduced structure.
It has relocated structure into an interpretive mechanism.

\paragraph{Unmodeled transformation behavior.}
A system is outside the scope when its essential identity behavior
depends on transformations not represented in $\Fset$ and not reducible to the
transformation families used here.
Such a case does not refute the lower bound.
It indicates that the transformation basis should be enlarged.
By monotonicity, such enlargement may preserve the bound or force additional regimes.

% -----------------------------------------------------------------------------
\subsection{Non-Goals and Weakened Objectives}
\label{subsec:limits_nongoals}
% -----------------------------------------------------------------------------

Several questions remain outside this paper.
This paper constructs no causal or normative extensions.
It takes no position on which causal account is correct or which normative interpretation should govern.
It offers no domain-specific modeling guidance.
It does not claim that all useful ontologies should be neutral substrates.
It does not claim that hidden regimes are always poor engineering.
Its claim is only that hidden regimes are incompatible with the specific objective analyzed here:
stable substrate-level reference under persistent disagreement.

If stability is weakened to permit context-dependent regime membership,
some distinctions may be represented as role-relative or parameter-dependent refinements of fewer base types.
That may be a reasonable engineering choice in coordinated settings.
But then the structural burden moves into coordination, versioning, governance,
or application logic whose admissibility is not fixed at the substrate level.
The lower bound proved here does not apply under that weakened objective.

% -----------------------------------------------------------------------------
\subsection{Boundary with Accountable Records}
\label{subsec:limits_boundary_accountable_records}
% -----------------------------------------------------------------------------

This paper concerns identity.
It derives the lower-bound regime structure required for stable reference.
It does not provide a conformance checker for deployed record systems.
It does not specify how schemas must declare bases, logs, or discriminator surfaces.
It does not decide whether an implementation has disclosed every identity-relevant surface.

Those questions are deferred to the companion \emph{Accountable Records} paper.
That paper is intended to develop accountable records, basis declaration,
transformation logs, and conformance checks.
The present paper supplies the lower-bound identity core that such a conformance analysis requires.
It establishes part of the identity cost of the neutral-substrate constraint,
while \emph{Accountable Records} studies how that cost
can be made checkable in deployed record systems.

% -----------------------------------------------------------------------------
\subsection{Boundary of the Claim}
\label{subsec:limits_boundary_claim}
% -----------------------------------------------------------------------------

Within the stated scope, the conclusion is unavoidable.
A neutral accountability substrate that supports the inventory of Section~\ref{sec:inventory},
preserves stable reference under the transformation basis $\Fset$,
and excludes hidden regimes must realize
at least the nine identity regimes in the lower-bound core.
Outside that scope, different optimization objectives may yield different structures.
The result is best understood as a constraint theorem:
for neutral substrate-level reference under persistent disagreement,
at least this much identity structure is required.
